\clearpage
\section*{October 2, 2026: A splitting cost from each historical parent's site}
\textit{Agent-drafted research entry.}

The model simulates only historical parents, so a splitting cost that depends on
the site needs only historical sites. Those parents were first filed in
2019--2022 and their zoning lots have had years to be recorded, so the 180-day
window, and the rush before 421-a expired, need not limit them. The tax-lot
area variant had made splitting costlier on larger lots, which looked like
project size rather than the site. This entry repeats it with the recorded
zoning lot.

\paragraph{Branch.} \path{tax_lot_splitting} now carries the jump-only burden
of October 1. The lot-group models and the lot-group bootstrap were ported to
it unchanged in substance.

\paragraph{A two-year window.} Of historical parents with a ZONE document from
two years before to three years after the first filing, about 30 percent are
first recorded more than 180 days after it, 13 percent more than a year after,
and 5 percent more than two years after. \path{build_parent_site_characteristics}
now computes the site measures over two windows after the first filing: 180
days, the same in both periods (the established columns, reproduced exactly),
and two years (\texttt{site\_lots\_two\_year},
\texttt{site\_lot\_area\_two\_year\_sqft} and companions), complete for every
2019--2022 parent and incomplete for most recent ones. Mergers and subdivisions
are undone, and demolitions counted, over the same window. Among the 588
historical parents of the comparison, 52 percent sit on a site of several lots
over two years, against 41 percent at 180 days and 26 percent by starting lots.
Site area remains tied to size: its correlation with log units is 0.65,
against 0.57 for tax-lot area.

\paragraph{Variants.} Two pooled variants of the main model scale each
historical parent's mean splitting cost by its site: by land,
$\sigma(A/\tilde A)^{-\beta}$ with $\tilde A$ the median historical site, or by
several lots, $\sigma e^{-\beta}$ on a site of two or more. The fitted cells and
the likelihood are those of the main model.

\begin{center}
\small
\begin{tabular}{lrrrr}
\hline
Model & $\beta$ & Log-likelihood & $p$ & Units lost \\
\hline
Main model & & $-584.5$ & & 761 \\
Tax-lot area & $-2$ & $-582.9$ & 0.07 & 1,697 \\
Site area, two years & $-3$ & $-583.1$ & 0.09 & 1,662 \\
Several site lots, two years & 4 & $-583.5$ & 0.15 & 1,406 \\
\hline
\end{tabular}
\end{center}

\paragraph{Findings.}
\begin{enumerate}
\item Site area behaves like tax-lot area: splitting is costlier on larger
sites, with $\beta$ and the spread $s = 4$ at the edges of their grids. The
correct site measure does not change the reading that land stands in for
project size.
\item Several site lots point the intuitive way: splitting is far cheaper on a
multi-lot site ($e^{-4}$, mean cost 0.046 against 2.5 on one lot). The fit
improves by one log point, not significantly, and the profile is flat: $\beta$
from $-1$ to 5 lies within 1.03 log points of the best.
\item Along those flat profiles units lost run from 761 to 1,756. Making
splitting depend on the site, in either form, raises units lost, because
parents that cannot split cheaply shrink instead.
\item Pooled cells carry little information about this: the site varies only
on the historical side, and the recent outcomes are counted by size and
building count alone. The lot-group fits, which also count where recent
splitters are, improve on their reference by likelihood ratios of 32--34, but
use the 180-day measure. The informative version therefore awaits recent
parents' complete windows.
\end{enumerate}

The main model is kept. The site variants join the alternatives behind the
range of units lost.
